\documentclass[a4paper,12pt]{article}
\usepackage[T2A]{fontenc}
\usepackage[utf8]{inputenc}
\usepackage[russian]{babel}
\usepackage{geometry}
\usepackage{listings}
\usepackage{longtable}
\geometry{margin=2.5cm}
\lstset{basicstyle=\ttfamily\small,breaklines=true,columns=fullflexible}

\begin{document}

\begin{titlepage}
\centering
{\large Национальный исследовательский университет ИТМО\par}
{\large Факультет программной инженерии и компьютерной техники\par}
\vspace{4cm}
{\large Отчёт к практическому заданию №1\par}
{\large по дисциплине «Системное программное обеспечение 2»\par}
\vfill
\begin{flushright}
Студент: Горинов Даниил Андреевич\\
Группа: P4116
\end{flushright}
\vfill
{Санкт-Петербург}\par
{2025/2026 учебный год\par}
\end{titlepage}

\clearpage
\pagenumbering{arabic}

\section{Цели}
Реализовать алгоритмы планирования задач в многопоточной среде на базе комплекса программ, разработанного в первом семестре, и показать различия в ходе выполнения тестовой программы при разных дисциплинах планирования.

Для варианта 4 реализованы алгоритмы \texttt{FCFS} и \texttt{SPN}. Критерий сравнения для чётного варианта — минимальное среднее время ожидания.

\section{Задачи}
\begin{enumerate}
\item Изучить работу таймера и прерываний в целевой виртуальной машине.
\item Добавить сохранение и восстановление контекста CPU для логических потоков.
\item Реализовать планировщики \texttt{FCFS} и \texttt{SPN}.
\item Подготовить одинаковую вычислительную нагрузку для двух алгоритмов.
\item Запустить тестовую программу через RemoteTasks и собрать метрики.
\item Сравнить среднее время ожидания, среднее время оборота и контрольные суммы потоков.
\end{enumerate}

\section{Описание работы}
Тестовая программа находится в \texttt{tests/scheduler\_demo.asm}. Она создаёт шесть логических потоков с одинаковым набором времён поступления и длительностей для двух запусков планировщика:
\begin{lstlisting}
arrivals = 0 3 6 9 12 15
bursts   = 4 4 7 7 4 4
\end{lstlisting}

Все длительности входят в диапазон варианта \texttt{3..7}, средний интервал поступления равен \texttt{3}.

Запуск выполняется командой:
\begin{lstlisting}
make -C SPO6 remote-demo
\end{lstlisting}

Команда собирает проект, копирует \texttt{tests/scheduler\_demo.asm} в \texttt{results/scheduler\_demo.asm}, запускает программу через \texttt{Portable.RemoteTasks.Manager.exe} и проверяет \texttt{results/scheduler\_demo.stdout.txt} скриптом \texttt{tools/check\_scheduler\_output.py}.

Основные артефакты:
\begin{itemize}
\item \texttt{spo6.target.pdsl} — описание целевой VM;
\item \texttt{devices.xml} — подключение \texttt{SimplePic} и \texttt{SimpleClock};
\item \texttt{tests/scheduler\_demo.asm} — демонстрация планировщиков;
\item \texttt{tests/timer\_probe.asm} — минимальная проверка таймера;
\item \texttt{results/scheduler\_demo.stdout.txt} — метрики выполнения.
\end{itemize}

\section{Аспекты реализации}
Таймер реализован через внешние устройства RemoteTasks. \texttt{SimpleClock} обновляет счётчик \texttt{Cycles} и формирует сигнал \texttt{on-cycles}; \texttt{SimplePic} принимает этот сигнал как прерывание 2, сохраняет адрес прерванной инструкции в \texttt{pic\_state} и передаёт управление обработчику из таблицы \texttt{pic\_handlers}.

В \texttt{spo6.target.pdsl} добавлены инструкции \texttt{SET\_CYCLES\_HANDLER}, \texttt{SET\_PERIOD}, \texttt{EI}, \texttt{DI}, \texttt{irq\_enter} и \texttt{iret}. Первая инструкция обработчика сохраняет \texttt{sp}, \texttt{bp}, \texttt{dbp} и \texttt{csp} в служебные ячейки \texttt{irq\_*}, затем переключает машину на kernel-context. \texttt{iret} восстанавливает сохранённые регистры и возвращает управление в поток.

Контекст каждого логического потока хранится в PCB. В PCB входят сохранённые регистры, время поступления, длительность, оставшееся время, состояние, время старта и завершения, а также контрольная сумма вычислений. Поверх PCB реализованы процедуры \texttt{create\_thread(entry\_ip, stack\_top, task\_end)} и \texttt{on\_thread\_interrupt(handler\_ip)}. Первая создаёт поток, вторая регистрирует пользовательский обработчик события перед возвратом из прерывания.

\texttt{FCFS} использует FIFO-очередь готовых задач. \texttt{SPN} при освобождении CPU выбирает готовую задачу с минимальной длительностью; при равенстве длительностей сохраняется порядок поступления. Оба алгоритма невытесняющие: таймер вызывает обработчик регулярно, но смена выполняемого потока происходит только в корректной точке завершения задачи или при простое CPU.

\section{Результаты}
Последний зафиксированный RemoteTasks-прогон:
\begin{lstlisting}
assemble = 231b82d3-5319-4f72-9a15-704ca2827dd4
run      = d561c55b-dc04-4d59-9c71-55c2d3d80530
\end{lstlisting}

\begin{center}
\begin{tabular}{|l|r|r|}
\hline
Метрика & FCFS & SPN \\
\hline
Среднее время ожидания \texttt{AW} & 5 & 4 \\
Среднее время оборота \texttt{AT} & 10 & 9 \\
Число dispatch-событий \texttt{D} & 6 & 6 \\
Число вызовов обработчика \texttt{I} & 30 & 30 \\
Число пользовательских хуков \texttt{H} & 30 & 30 \\
\hline
\end{tabular}
\end{center}

Контрольные суммы совпали для обоих запусков:
\begin{lstlisting}
406 806 2121 2821 2006 2406
\end{lstlisting}

Совпадение контрольных сумм показывает, что различается только порядок планирования, а не вычисления внутри потоков.

\section{Выводы}
Таймер, прерывания и переключение контекста реализованы на уровне целевой VM и runtime-обработчика. Тестовая программа показала различие дисциплин планирования на одной нагрузке: \texttt{SPN} уменьшил среднее время ожидания с \texttt{5} до \texttt{4} и среднее время оборота с \texttt{10} до \texttt{9}, но увеличил ожидание длинной четвёртой задачи. Для чётного варианта по критерию среднего времени ожидания предпочтителен \texttt{SPN}.

\end{document}
